% =================================================================
%  style/tikz-styles.tex
%  TikZ и наборы стилей для иллюстраций матанализа.
%  pgfplots в эталоне не используется: графики рисуются
%  через \draw plot (\x,{...}) внутри обычного tikzpicture.
% =================================================================

\usepackage{tikz}
\usetikzlibrary{arrows.meta,patterns,calc}

\tikzset{
	% --- оси и сетка ---
	ax/.style        = {thick, -{Stealth[length=3mm]}},
	gridline/.style  = {step=0.5, gray!25, very thin},
	tickmark/.style  = {thick},
	% --- кривые ---
	curve/.style     = {thick, smooth, Blue},
	curvehl/.style   = {ultra thick, Purple},
	curveaux/.style  = {thick, dashed, black!50},
	% контур открытой области
	curvebd/.style  = {thick, densely dashed, black!50},
	% контур замкнутой области и фоновые кривые: сплошные, но не акцентные
	region/.style    = {thick, smooth, black!50},
	% серые штрихи: поле направлений и штриховка областей
	curvehatch/.style = {black!40, line width=0.4pt},
	% --- заливки-штриховки под кривой / прямоугольники сумм ---
	hatchBlue/.style    = {pattern=north east lines, pattern color=hlBlue, opacity=0.85},
	hatchPurple/.style    = {pattern=north east lines, pattern color=hlPurple,   opacity=0.85},
	hatchGray/.style = {pattern=north east lines, pattern color=black!60, opacity=0.25},
	% --- векторы и проекции ---
	vecarr/.style    = {->, very thick},
	vecaux/.style    = {->, dashed},
	pt/.style        = {circle, fill, inner sep=1.3pt},
	% выколотая точка: белый кружок с серым контуром
	ptopen/.style    = {circle, draw=black!50, fill=white,
	                    line width=0.6pt, inner sep=1.5pt},
	% --- подписи ---
	llbl/.style       = {font=\large},
	lbl/.style       = {font=\small},
	slbl/.style      = {font=\footnotesize},
}

% -----------------------------------------------------------------
% Прямой угол в точке (#1) со сторонами вдоль осей:
%     \rightangle{(P)}       — по умолчанию влево-вверх
% -----------------------------------------------------------------
\newcommand{\rightangle}[2][0.28]{%
	\draw #2 ++(-#1,0) -- ++(0,#1) -- ++(#1,0);%
}
